% TEMPLATE-MILC-v3.tex - plantilla maestra para todas las guias MILC v3
% Ver scripts/generadores/build-guia-milc.py para la lista completa de
% claves esperadas. El script valida que no queden placeholders sin
% reemplazar antes de escribir el archivo final.

\documentclass[11pt,letterpaper]{article}
\usepackage[margin=1.75cm]{geometry}
\usepackage[table]{xcolor}
\usepackage{fontspec}
\usepackage[spanish]{babel}
\usepackage{tikz}
\usepackage[most]{tcolorbox}
\usepackage{tabularx}
\usepackage{array}
\usepackage{fancyhdr}
\usepackage{titlesec}
\usepackage{ragged2e}
\usepackage{enumitem}
\usepackage{microtype}

\setmainfont{Helvetica Neue}
\setsansfont{Helvetica Neue}
\setlength{\parindent}{0pt}
\setlength{\parskip}{6pt}
\hyphenpenalty=200
\exhyphenpenalty=200
\pretolerance=400
\tolerance=2000
\setlength{\emergencystretch}{1em}
\renewcommand{\arraystretch}{1.25}
\setlist{leftmargin=5mm,itemsep=1mm,topsep=1mm}
\newcolumntype{Y}{>{\RaggedRight\arraybackslash}X}

\definecolor{milcMagenta}{HTML}{E6007E}
\definecolor{milcVino}{HTML}{5A0038}
\definecolor{milcTurquesa}{HTML}{0093A5}
\definecolor{milcVerde}{HTML}{58A923}
\definecolor{milcMostaza}{HTML}{E5B400}
\definecolor{milcOcre}{HTML}{B86600}
\definecolor{milcNegro}{HTML}{241816}
\definecolor{milcGris}{HTML}{F7F7F4}
\definecolor{milcLinea}{HTML}{E8E2DD}
\definecolor{milcGradePrimary}{HTML}{FF2D87}
\definecolor{milcGradeDark}{HTML}{9F1B56}
\definecolor{milcGradeSoft}{HTML}{FFE0EE}
\definecolor{milcGradeAccent}{HTML}{CFE7FF}
\definecolor{milcGradeText}{HTML}{FFFFFF}
\color{milcNegro}

\pagestyle{fancy}
\fancyhf{}
\lhead{\small Tecnología e Informática · Grado 10}
\rhead{\small Guía 15 · MILC v3}
\cfoot{\small\thepage}
\renewcommand{\headrulewidth}{0pt}

\newtcolorbox{titlebox}[1]{
  enhanced,colback=#1,colframe=#1,arc=7mm,boxrule=0pt,
  left=7mm,right=7mm,top=3mm,bottom=3mm,
  before skip=8pt,after skip=8pt,
  fontupper=\sffamily\bfseries\Large\color{white}
}

\newtcolorbox{softbox}[3]{
  enhanced,colback=#2,colframe=#1,borderline west={4pt}{0pt}{#1},
  arc=2mm,boxrule=.4pt,left=4mm,right=4mm,top=3mm,bottom=3mm,
  before skip=6pt,after skip=8pt,title={#3},coltitle=white,
  colbacktitle=#1,fonttitle=\sffamily\bfseries,fontupper=\RaggedRight,
  attach boxed title to top left={xshift=3mm,yshift=-2mm},
  boxed title style={arc=2mm, boxrule=0pt}
}

\newtcolorbox{drawbox}[2]{
  enhanced,colback=white,colframe=#1,arc=4mm,boxrule=1pt,
  left=5mm,right=5mm,top=4mm,bottom=4mm,before skip=8pt,after skip=8pt,
  title={#2},coltitle=white,colbacktitle=#1,fonttitle=\sffamily\bfseries,
  fontupper=\RaggedRight,
  attach boxed title to top left={xshift=4mm,yshift=-2mm},
  boxed title style={arc=3mm, boxrule=0pt}
}

\newtcolorbox{infoband}[1]{
  enhanced,colback=#1!8,colframe=#1!50,arc=2mm,boxrule=.5pt,
  left=4mm,right=4mm,top=2mm,bottom=2mm,before skip=4pt,after skip=4pt,
  fontupper=\small\RaggedRight
}

% Puente narrativo entre secciones (contrato editorial Conéctate).
% Da continuidad: "Acabas de X. Ahora vas a Y porque Z."
% Visualmente: banda izquierda en color del grado, texto en itálica pequeña.
\newtcolorbox{puentebox}{
  enhanced,colback=milcGradeSoft,colframe=milcGradeDark,
  borderline west={3pt}{0pt}{milcGradeDark},
  arc=0mm,boxrule=0pt,left=5mm,right=4mm,top=2.5mm,bottom=2.5mm,
  before skip=4pt,after skip=8pt,
  fontupper=\itshape\small\color{milcNegro}\RaggedRight
}
\newcommand{\puente}[1]{%
  \begin{puentebox}%
    \textcolor{milcGradeDark}{\textbf{\textrightarrow}}\hspace{2mm}#1%
  \end{puentebox}%
}

\newcommand{\linead}[1]{\noindent\color{milcLinea}\rule{#1}{0.5pt}\color{milcNegro}}
\newcommand{\respuesta}[1]{\par\vspace{2mm}\foreach \n in {1,...,#1}{\noindent\color{milcLinea}\rule{\linewidth}{0.5pt}\par\vspace{5mm}}\color{milcNegro}}
\newcommand{\checkbox}{\textcolor{milcGradeDark}{\fbox{\phantom{x}}}\hspace{5pt}}

\newcommand{\phasecard}[4]{
\begin{tcolorbox}[enhanced,colback=#3,colframe=#2,arc=3mm,boxrule=.5pt,height=3.05cm,valign=center,halign=center,left=1.5mm,right=1.5mm,top=2mm,bottom=2mm]
{\sffamily\bfseries\fontsize{22}{24}\selectfont\textcolor{#2}{#1}}\par
{\sffamily\bfseries\small #4}
\end{tcolorbox}
}

\begin{document}
\thispagestyle{empty}
\begin{tikzpicture}[remember picture,overlay]
  \fill[white] (current page.south west) rectangle (current page.north east);
  \path[fill=milcGradePrimary,rounded corners=16mm]
    ([xshift=.55cm,yshift=-.55cm]current page.north west) rectangle
    ([xshift=-.55cm,yshift=3.65cm]current page.south east);

  \node[anchor=north west,text=milcGradeText] at ([xshift=2.0cm,yshift=-1.85cm]current page.north west)
    {\sffamily\bfseries\fontsize{14}{16}\selectfont GRADO};
  \node[anchor=north west,text=milcGradeText,opacity=.82] at ([xshift=2.0cm,yshift=-2.75cm]current page.north west)
    {\sffamily\bfseries\fontsize{9}{11}\selectfont SERIE GUÍAS · TECNOLOGÍA E INFORMÁTICA};

  \begin{scope}[shift={([xshift=-3.05cm,yshift=-2.55cm]current page.north east)}]
    \fill[white,opacity=.80] (-.62,-.52) rectangle (.62,.52);
    \draw[milcGradeText,line width=1pt] (-.62,-.52) rectangle (.62,.52);
    \fill[milcGradeDark] (-.42,-.38) rectangle (-.22,.06);
    \fill[milcGradeAccent] (-.10,-.38) rectangle (.10,.24);
    \fill[milcGradePrimary!60!black] (.22,-.38) rectangle (.42,.38);
  \end{scope}

  \node[anchor=west,text=milcGradeText] at ([xshift=1.9cm,yshift=-12.85cm]current page.north west)
    {\sffamily\bfseries\fontsize{124}{124}\selectfont 10};
  \draw[milcGradeText,line width=1.7pt]
    ([xshift=4.52cm,yshift=-11.68cm]current page.north west) circle (.13cm);

  \node[anchor=west,text=milcGradeText,text width=8.7cm,align=left] at ([xshift=10.35cm,yshift=-11.6cm]current page.north west)
    {
     {\sffamily\bfseries\fontsize{19}{22}\selectfont Encuestas con phronesis ---\\Google Forms\\más IA}\\[4mm]
     {\sffamily\bfseries\fontsize{23}{26}\selectfont Guía 15-10-TIC}\\[2mm]
     {\sffamily\fontsize{13}{16}\selectfont Período 2 · Informes técnicos y comunicación profesional con IA}\\[5mm]
     {\sffamily\bfseries\fontsize{11}{13}\selectfont Institución Educativa Sor María Juliana}\\[1mm]
     {\sffamily\fontsize{10}{12}\selectfont Autor: Dr. Álvaro Cárdenas Orozco}
    };

  \path[fill=milcGradeSoft,rounded corners=7mm]
    ([xshift=1.35cm,yshift=.85cm]current page.south west) rectangle
    ([xshift=-1.35cm,yshift=4.25cm]current page.south east);
  \draw[milcGradeDark,line width=.65pt,rounded corners=7mm]
    ([xshift=1.35cm,yshift=.85cm]current page.south west) rectangle
    ([xshift=-1.35cm,yshift=4.25cm]current page.south east);
  \node[anchor=center,text=milcNegro,text width=17.0cm,align=center] at ([yshift=2.55cm]current page.south)
    {\sffamily\fontsize{10}{12}\selectfont Metodología MILC · Apertura ancestral · Pensamiento computacional · Triángulo Dussel-Estoicismo-Floridi\\
    \textcolor{milcGradeDark}{\bfseries Producto:} Encuesta diseñada en Google Forms con al menos 10 preguntas (mezcla de tipos: opción única, opción múltiple, escala Likert 1-5, abiertas, sí/no) sobre un tema escolar real (vida del recreo, consumo de tienda escolar, tiempo de pantalla, transporte al colegio, gustos musicales del grupo) + revisión con IA del cuestionario completo para detectar preguntas sesgadas, ambiguas o capciosas con bitácora de los ajustes + plan de aplicación a 20-30 personas con cronograma y consentimiento informado};
\end{tikzpicture}
\mbox{}
\newpage

\section*{Apertura: Encuestas con phronesis --- Google Forms más IA para preguntar bien}

\begin{softbox}{milcGradeDark}{milcGradeSoft}{Saber ancestral}
En cualquier barrio popular del Valle del Cauca, en cualquier vereda del Cauca o del Quindío, hay una figura silenciosa que cualquier vecino mayor reconoce con una mezcla de respeto y prudencia: \textbf{la persona que sabe preguntar sin parecer entrometida}. No es chismoso barato (\emph{``y a usted qué le importa''}); es el que averigua con método. Su técnica es ancestral y precisa: (1) \textbf{Pregunta indirecta primero}: empieza por algo aparentemente banal (\emph{``¿qué tal el tiempo hoy, doña Rosa?''}). (2) \textbf{Observa la cara}: ve si la persona se relaja, se incomoda, evade o se abre. (3) \textbf{Profundiza con calma}: si ve apertura, sigue. Si ve resistencia, retrocede. (4) \textbf{Nunca pregunta capciosa}: no induce respuestas con preguntas tendenciosas. (5) \textbf{Respeta el silencio}: deja espacio para que la otra persona piense. Esa persona logra enterarse de lo que pasa en el barrio sin que nadie se sienta interrogado. La sabiduría ancestral es phronesis pura: \textbf{preguntar bien es oficio antes que técnica}. Esa misma phronesis es la que pide una buena encuesta profesional: orden, neutralidad, respeto por el encuestado.
\end{softbox}

\begin{softbox}{milcTurquesa}{milcGris}{Saber contemporáneo}
Una \textbf{encuesta profesional} es instrumento de recolección de datos que requiere diseño cuidadoso. La regla profesional: \textbf{las respuestas dependen de cómo preguntas, no solo de a quién preguntas}. \textbf{Google Forms} (gratis con cuenta Gmail) es la herramienta más usada para encuestas digitales. Soporta varios \textbf{tipos de preguntas}: (1) \textbf{Opción única (multiple choice)}: el encuestado elige 1 de varias opciones. Útil para variables categóricas (sexo, transporte usado, grado). (2) \textbf{Opción múltiple (checkbox)}: el encuestado elige varias. Útil para opciones que no son excluyentes. (3) \textbf{Escala Likert 1-5}: medición de actitudes (1=muy en desacuerdo, 5=muy de acuerdo). (4) \textbf{Respuesta abierta corta}: para datos específicos (edad, número). (5) \textbf{Respuesta abierta larga}: para comentarios cualitativos. (6) \textbf{Cuadrícula}: para evaluar varios items con la misma escala. La regla profesional fundamental: \textbf{toda pregunta debe pasar el filtro de la IA antes de aplicarse}. ChatGPT, Gemini o Claude pueden detectar 5 problemas comunes: preguntas sesgadas, ambiguas, capciosas, doble cabeza, asumir conocimiento previo. Pedirle a la IA \emph{``analiza este cuestionario para detectar preguntas problemáticas''} es práctica profesional moderna.
\end{softbox}

\begin{softbox}{milcMostaza}{milcGris}{Pregunta-puente}
\textit{¿Qué sabía el que pregunta bien del barrio al respetar el silencio y observar la cara, que el novato de Forms olvida cuando hace preguntas tendenciosas como \emph{``¿no le parece que el recreo es muy corto?''}? ¿Y por qué pasar el cuestionario por IA antes de aplicarlo evita 80 por ciento de los errores comunes?}
\end{softbox}

\begin{softbox}{milcVerde}{milcGris}{Saber y saber hacer}
Al terminar podrás: (1) \textbf{identificar} los 6 tipos de preguntas de Google Forms y cuándo usar cada uno; (2) \textbf{analizar} un cuestionario tuyo con asistencia de IA para detectar sesgos, ambigüedades y preguntas capciosas; (3) \textbf{aplicar} la phronesis del preguntón ancestral al diseño de preguntas neutras y útiles; (4) \textbf{evaluar} el cuestionario final ajustando según el análisis con IA.
\end{softbox}

\small
\begin{tabularx}{\textwidth}{>{\bfseries}p{3.0cm}Y}
\rowcolor{milcGradePrimary!12}Autor & Dr. Álvaro Cárdenas Orozco\\
DBA & Comunicación profesional con asistencia de IA y herramientas ofimáticas gratuitas (MEN, T\&I)\\
Referentes & Markdown como estándar abierto · Google Workspace · Estoicismo (claridad)\\
Producto final & Encuesta diseñada en Google Forms con al menos 10 preguntas (mezcla de tipos: opción única, opción múltiple, escala Likert 1-5, abiertas, sí/no) sobre un tema escolar real (vida del recreo, consumo de tienda escolar, tiempo de pantalla, transporte al colegio, gustos musicales del grupo) + revisión con IA del cuestionario completo para detectar preguntas sesgadas, ambiguas o capciosas con bitácora de los ajustes + plan de aplicación a 20-30 personas con cronograma y consentimiento informado\\
\end{tabularx}
\normalsize

\puente{Antes de empezar, mira el viaje completo. En las sesiones 1-4 trabajaste textos profesionales escritos. Hoy aprendes a \emph{recolectar datos propios} con encuestas. Las sesiones 6-9 usarán los datos que recolectes hoy para análisis con IA y luego informe comercial. Sin buenos datos hoy, los análisis de mañana serán falsos.}

\begin{titlebox}{milcGradeDark}
Ruta MILC: así vamos a aprender
\end{titlebox}
\begin{tabularx}{\textwidth}{>{\centering\arraybackslash}X>{\centering\arraybackslash}X>{\centering\arraybackslash}X>{\centering\arraybackslash}X}
\phasecard{1}{milcVerde}{milcGris}{Escuta\\[-1mm]\small Observo, escucho y pregunto.} &
\phasecard{2}{milcTurquesa}{milcGris}{Sistematización\\[-1mm]\small Pensamiento computacional.} &
\phasecard{3}{milcMagenta}{milcGris}{Praxis\\[-1mm]\small Construyo y aplico.} &
\phasecard{4}{milcVino}{milcGris}{Evaluación\\[-1mm]\small Reflexión liberadora.}
\end{tabularx}


\puente{Antes de teorizar, vas a hacer un \emph{experimento del entrevistador novato}: imagina que quieres saber si los compañeros del grupo están de acuerdo con cambiar el horario del recreo. Escribe 3 preguntas que harías. Después, mírate críticamente: ¿alguna de tus preguntas induce la respuesta que tú esperas?}

\begin{titlebox}{milcVerde}
Fase 1 · Escuta
\end{titlebox}

\begin{softbox}{milcVerde}{milcGris}{Escena para escuchar}
\textbf{Actividad 1 · IDENTIFICA --- 3 preguntas críticas} (15 min · individual). Imagina que quieres saber qué piensa tu grupo sobre el horario actual del recreo (15 minutos). Sin investigar nada todavía, escribe 3 preguntas que harías. Después léelas con ojo crítico y verifica: (1) ¿alguna induce la respuesta que tú esperas? (\emph{``¿no le parece que 15 minutos es muy poco?''}). (2) ¿alguna asume conocimiento que el encuestado puede no tener? (3) ¿alguna tiene 2 preguntas dentro de una sola (doble cabeza)? Esa autocrítica te entrena el ojo para detectar sesgos.
\end{softbox}

\checkbox Escribí 3 preguntas iniciales sobre el recreo.\\
\checkbox Las revisé críticamente.\\
\checkbox Identifiqué si tienen sesgos, ambigüedades o doble cabeza.

\begin{infoband}{milcVerde}
\textbf{Lo que va al cuaderno (Actividad 1 · IDENTIFICA).} Título: «Actividad 1 --- 3 preguntas críticas». \textbf{Formato:} 3 preguntas iniciales + autocrítica de cada una. \textbf{Sabes que terminaste cuando} ves que tus propias preguntas iniciales tienen sesgos.
\end{infoband}

\puente{Ya tienes tus 3 preguntas críticas. Ahora vas a entender la \textbf{anatomía de las encuestas profesionales}: los 6 tipos de pregunta, los 5 errores comunes, el filtro de IA.}

\begin{titlebox}{milcTurquesa}
Fase 2 · Sistematización con pensamiento computacional
\end{titlebox}

\begin{softbox}{milcTurquesa}{milcGris}{Cuatro pilares aplicados al tema}
Tu autocrítica te mostró que los sesgos son fáciles de cometer. Ahora necesitas la \textbf{anatomía profesional} de las encuestas y el filtro de IA.
\end{softbox}

\small
\begin{tabularx}{\textwidth}{>{\bfseries\RaggedRight\arraybackslash}p{3.6cm}Y}
\rowcolor{milcTurquesa!90}\color{white}Pilar & \color{white}\bfseries Aplicación al tema\\
1. Descomposición & Los \textbf{6 tipos de preguntas de Google Forms} se descomponen así: (1) \textbf{Opción única}: \emph{``¿Cuál es tu medio de transporte al colegio?''} (a) Pie (b) Bus (c) Bicicleta (d) Moto (e) Otro. (2) \textbf{Opción múltiple}: \emph{``¿Qué actividades haces durante el recreo?''} (puede marcar varias). (3) \textbf{Escala Likert 1-5}: \emph{``Estoy satisfecho con la duración del recreo''} (1=muy en desacuerdo, 5=muy de acuerdo). (4) \textbf{Respuesta abierta corta}: \emph{``¿Cuántos minutos durmiste anoche?''} (5) \textbf{Respuesta abierta larga}: \emph{``Si pudieras cambiar algo del recreo, ¿qué sería y por qué?''}. (6) \textbf{Cuadrícula}: tabla que evalúa varios items con la misma escala (\emph{``Califica de 1 a 5 cada uno: cafetería, baños, recreo, biblioteca''}).\\
2. Reconocimiento de patrones & Los \textbf{5 errores comunes} en encuestas y cómo evitarlos: (a) \textbf{Pregunta sesgada}: \emph{``¿No le parece que el recreo debería ser más largo?''} (induce sí). \emph{Mejor:} \emph{``¿La duración actual del recreo es adecuada para ti?''}. (b) \textbf{Pregunta ambigua}: \emph{``¿Te gusta la cafetería?''}. ¿Te gusta qué de la cafetería? \emph{Mejor:} preguntar por aspectos específicos (precios, sabor, opciones, atención). (c) \textbf{Doble cabeza}: \emph{``¿La biblioteca y la sala de cómputo te parecen útiles?''} son 2 preguntas. \emph{Mejor:} separarlas. (d) \textbf{Asumir conocimiento previo}: \emph{``¿Apoyas el nuevo sistema MILC?''} sin explicar qué es. (e) \textbf{Opciones de respuesta incompletas}: si preguntas \emph{``¿qué transporte usas?''} debes incluir \emph{otro} para casos no previstos.\\
3. Abstracción & El \textbf{filtro de IA} es práctica profesional moderna. El prompt típico: \emph{``Actúa como experto en metodología de encuestas. Aquí está mi cuestionario sobre [tema] aplicado a [audiencia]. Analiza cada pregunta y detecta: (1) sesgos (preguntas que inducen respuesta), (2) ambigüedades, (3) preguntas de doble cabeza, (4) asunciones de conocimiento previo, (5) opciones de respuesta incompletas. Sugiere ajustes específicos''}. Pasa el resultado por ChatGPT, Gemini o Claude. La IA suele detectar el 80 por ciento de los problemas comunes.\\
4. Algoritmo conceptual & Algoritmo: (1) elegir tema real del colegio. (2) escribir 10 preguntas iniciales con mezcla de tipos. (3) revisarlas a mano con la lista de 5 errores. (4) pasarlas por IA para análisis completo. (5) ajustar según hallazgos. (6) crear encuesta en Forms. (7) preparar plan de aplicación con consentimiento informado.\\
\end{tabularx}
\normalsize

\begin{drawbox}{milcTurquesa}{Anatomía de la encuesta profesional}
\textbf{6 tipos:} única / múltiple / Likert / abierta corta / abierta larga / cuadrícula. \\[1mm] \textbf{5 errores:} sesgo / ambigüedad / doble cabeza / asunción / opciones incompletas. \\[1mm] \textbf{Filtro IA:} análisis con prompt específico antes de aplicar. \\[1mm] \textbf{Consentimiento:} declarar propósito y anonimato.
\end{drawbox}

\begin{softbox}{milcMostaza}{milcGris}{Errores comunes y su corrección}
(1) Pregunta sesgada que induce respuesta. (2) Pregunta ambigua sin especificar aspecto. (3) Doble cabeza (2 preguntas en 1). (4) Asumir conocimiento que el encuestado no tiene. (5) Opciones de respuesta incompletas (sin \emph{otro}). (6) Encuesta demasiado larga (más de 15 preguntas cansa). (7) Sin consentimiento informado al inicio.
\end{softbox}

\begin{infoband}{milcTurquesa}
\textbf{Lo que va al cuaderno (Actividad 2 · ANALIZA).} Título: «Actividad 2 --- Anatomía encuesta». \textbf{Formato:} (a) los 6 tipos con ejemplo; (b) los 5 errores con corrección; (c) prompt de filtro IA. \textbf{Sabes que terminaste cuando} sabes detectar los 5 errores en preguntas ajenas.
\end{infoband}

\puente{Tienes el método. Toca aplicarlo: diseñas encuesta en Forms con 10 preguntas, la pasas por IA para análisis, ajustas según hallazgos, preparas plan de aplicación a 20-30 personas.}

\begin{titlebox}{milcMagenta}
Fase 3 · Praxis con pensamiento computacional
\end{titlebox}

\begin{softbox}{milcMagenta}{milcGris}{Construyendo el producto con los cuatro pilares}
\textbf{Actividad 3 · APLICA + EVALÚA --- Encuesta Forms + análisis IA + ajustes + plan} (40 min · individual con dispositivo). Hora de aplicar. Diseñas encuesta de 10 preguntas, la pasas por IA, ajustas, creas en Forms, preparas plan de aplicación.
\end{softbox}

\small
\begin{tabularx}{\textwidth}{>{\bfseries\RaggedRight\arraybackslash}p{3.6cm}Y}
\rowcolor{milcMagenta!90}\color{white}Pilar & \color{white}\bfseries Aplicación al producto\\
Descomposición & Tu encuesta se construye en 6 pasos: (1) elegir tema real (recreo, cafetería, transporte, tiempo de pantalla, hábitos de estudio). (2) escribir 10 preguntas mezclando 4-5 tipos distintos. (3) pasar las 10 preguntas por IA usando el prompt de filtro. (4) ajustar según el análisis (corregir sesgos, especificar ambigüedades, separar dobles cabezas). (5) crear la encuesta en Google Forms agregando consentimiento informado al inicio. (6) preparar plan de aplicación con cronograma y lista de 20-30 personas a encuestar.\\
Patrones & El \textbf{consentimiento informado al inicio} de la encuesta tiene 4 elementos: (a) propósito del estudio en 1-2 líneas. (b) anonimato declarado o nombre opcional. (c) tiempo estimado (\emph{``5 minutos''}). (d) opción de no participar. Sin estos 4 elementos, la encuesta es éticamente cuestionable.\\
Abstracción & La bitácora de ajustes tiene 3 columnas: (a) pregunta original. (b) error detectado por IA. (c) pregunta ajustada. Al menos 3-5 filas que demuestren cómo la IA te ayudó a mejorar.\\
Algoritmo de armado & Algoritmo: (1) tema. (2) 10 preguntas iniciales. (3) filtro IA. (4) ajustes con bitácora. (5) Forms con consentimiento. (6) plan de aplicación. (7) verificar tiempo total (debería tomar 5-7 min responder).\\
\end{tabularx}
\normalsize

\begin{drawbox}{milcMagenta}{Checklist antes de entregar}
\checkbox 10 preguntas mezclando 4+ tipos. \\[1mm] \checkbox Filtro de IA aplicado al cuestionario completo. \\[1mm] \checkbox Bitácora con 3+ ajustes documentados. \\[1mm] \checkbox Encuesta creada en Google Forms. \\[1mm] \checkbox Consentimiento informado al inicio. \\[1mm] \checkbox Plan de aplicación a 20-30 personas. \\[1mm] \checkbox Tiempo total de respuesta verificado.
\end{drawbox}

\begin{softbox}{milcMagenta}{milcGris}{Plantilla de trabajo}
\textbf{Tema.} \textit{Escolar real.} \\[1mm] \textbf{10 preguntas.} \textit{Mezcla de 4+ tipos.} \\[1mm] \textbf{Filtro IA.} \textit{Prompt de análisis.} \\[1mm] \textbf{Bitácora.} \textit{3+ ajustes documentados.} \\[1mm] \textbf{Forms.} \textit{Encuesta con consentimiento.} \\[1mm] \textbf{Plan.} \textit{20-30 personas + cronograma.}
\end{softbox}

\begin{infoband}{milcMagenta}
\textbf{Lo que va al cuaderno (Actividad 3 · APLICA + EVALÚA).} Título: «Actividad 3 --- Mi encuesta profesional». \textbf{Formato:} (a) link a la encuesta Forms; (b) bitácora de ajustes; (c) plan de aplicación. \textbf{Sabes que terminaste cuando} la encuesta está lista para enviar el lunes.
\end{infoband}

\puente{Lo que sigue resume \emph{qué} entregas hoy: encuesta Forms + análisis IA + bitácora ajustes + plan de aplicación. El criterio principal: que las preguntas sean neutras y útiles para tomar decisión real.}

\begin{titlebox}{milcMostaza}
Producto de la sesión
\end{titlebox}
\textbf{Encuesta Forms + análisis IA + bitácora + plan aplicación}.
\begin{softbox}{milcMostaza}{milcGris}{Criterios de calidad}
(1) 10 preguntas con 4+ tipos. (2) Filtro de IA aplicado. (3) Bitácora de ajustes documentada. (4) Consentimiento informado al inicio. (5) Plan de aplicación a 20-30 personas. (6) Tiempo de respuesta razonable (5-7 min).
\end{softbox}

\puente{Antes de cerrar, mira la encuesta desde las cinco dimensiones humanas. Preguntar bien es respeto por el encuestado; preguntar mal es manipulación disfrazada de investigación.}

\begin{titlebox}{milcVino}
Fase 4 · Evaluación liberadora — cinco dimensiones
\end{titlebox}

\begin{softbox}{milcVino}{milcGris}{Sentido de la evaluación}
Evaluar no es solo revisar una nota. Es reconocer qué aprendí, qué cambió en mí, qué emociones manejé, cómo me relaciono con la ciudad y la comunidad, y qué le dejaré a quien viene detrás.
\end{softbox}

\textbf{Mi reflexión liberadora en cinco dimensiones.} Completa cada frase iniciada:

\small
\begin{tabularx}{\textwidth}{>{\bfseries\RaggedRight\arraybackslash}p{3.4cm}Y>{\RaggedRight\arraybackslash}p{4.6cm}}
\rowcolor{milcVino}\color{white}Dimensión & \color{white}\bfseries Mi respuesta (frame starter) & \color{white}\bfseries Algo que descubrí\\
1. Personal & Lo que más cambió en mí fue \linead{6.5cm} & \linead{4.2cm}\\
2. Emocional & La emoción más fuerte fue \linead{2.5cm} y la regulé \linead{3cm} & \linead{4.2cm}\\
3. Ciudadana & En democracia esto importa porque \linead{6cm} & \linead{4.2cm}\\
4. Local (Cartago) & En mi barrio sería útil para \linead{6.5cm} & \linead{4.2cm}\\
5. Intergeneracional & A mi abuela/abuelo le diría \linead{6.5cm} & \linead{4.2cm}\\
\end{tabularx}
\normalsize

\begin{infoband}{milcVino}
\textbf{Para la próxima sustentación.} Vuelve a esta tabla en seis meses. Verás que las respuestas cambian — no porque lo aprendido cambie, sino porque tú cambias. La evaluación liberadora es una conversación contigo a través del tiempo.
\end{infoband}


\puente{Y para terminar, tres voces sobre la palabra que pregunta. Dussel sobre las encuestas que dignifican al encuestado, Séneca sobre la neutralidad como virtud, Floridi sobre el dato bien recolectado como ética informacional.}

\begin{titlebox}{milcGradeDark}
Triángulo de pensamiento · cierre filosófico
\end{titlebox}

\begin{softbox}{milcVino}{milcGris}{Dussel · lente del nosotros}
\textit{``Una encuesta neutra dignifica al encuestado; una tendenciosa lo manipula bajo disfraz de investigación.''}

Cuando una encuesta tiene preguntas tendenciosas, no investiga: confirma lo que el encuestador ya quería oír. La filosofía de la liberación pide lo opuesto: \textbf{respetar al encuestado dejándole responder libremente}. La phronesis del preguntón ancestral hacía eso por instinto; la encuesta profesional moderna lo hace por método.

\textbf{¿Mi encuesta deja libertad al encuestado para responder honestamente, o lo induce a confirmar mi hipótesis?}
\end{softbox}
\linead{\linewidth}\\[1mm]
\linead{\linewidth}

\begin{softbox}{milcMostaza}{milcGris}{Estoicismo · Séneca · cuidado interior}
\textit{``La neutralidad en la pregunta es virtud del que busca verdad; la tendenciosidad es vicio del que ya tiene conclusión.''}

Es tentador formular preguntas que confirmen lo que tú piensas. La sabiduría estoica recomienda lo opuesto: \emph{si ya sabes la respuesta, no preguntes}. Si preguntas, prepárate para que el dato te sorprenda. La phronesis del oficio consiste en \textbf{preguntar para aprender, no para confirmar}.

\textbf{¿Estoy preguntando para aprender, o para que confirmen lo que ya pienso?}
\end{softbox}
\linead{\linewidth}\\[1mm]
\linead{\linewidth}

\begin{softbox}{milcTurquesa}{milcGris}{Floridi · lente de la infoesfera}
\textit{``El dato bien recolectado es ética informacional en la era de la sobreinformación.''}

En la era de la información, hay datos por todas partes pero pocos están bien recolectados. La ética profesional pide que cuando produces datos propios, los produzcas con rigor: pregunta neutra, consentimiento informado, anonimato si aplica, filtro de IA contra sesgos. Tu encuesta de hoy te entrena en esa práctica.

\textbf{¿Mis datos están siendo recolectados con rigor profesional, o solo replico la mala práctica que circula?}
\end{softbox}
\linead{\linewidth}\\[1mm]
\linead{\linewidth}

\begin{titlebox}{milcVino}
Compromiso final
\end{titlebox}

\small
\begin{tabularx}{\textwidth}{>{\RaggedRight\arraybackslash}p{3.0cm}YYY}
\rowcolor{milcVino}\color{white}\bfseries Criterio & \color{white}\bfseries Lo logré & \color{white}\bfseries En proceso & \color{white}\bfseries Necesito apoyo\\
Comprensión del tema & Explico ideas centrales con mis palabras. & Reconozco ideas pero falta precisión. & Necesito repasar conceptos base.\\
Pensamiento computacional & Apliqué los 4 pilares al diseñar mi producto. & Apliqué algunos pilares. & Necesito releer la fase 2.\\
Aplicación y producto & Mi producto cumple los criterios y se puede revisar. & Producto funcional con ajustes pendientes. & Aún en construcción.\\
Reflexión 5 dimensiones & Respondí cada dimensión con honestidad. & Respondí algunas con detalle. & Mis respuestas son muy generales.\\
Triángulo de pensamiento & Conecté las 3 voces con mi trabajo real. & Conecté al menos una. & No relaciono las voces con mi proceso.\\
\end{tabularx}
\normalsize

\textbf{Hoy entendí que:}\\[1mm]
\linead{\linewidth}\\[1mm]
\linead{\linewidth}

\textbf{Mi compromiso para la próxima sesión es:}\\[1mm]
\linead{\linewidth}\\[1mm]
\linead{\linewidth}

\begin{softbox}{milcMostaza}{milcGris}{Cita de despedida · Marco Aurelio}
\textit{``El obstáculo en el camino se vuelve el camino. Lo que estorba al camino se vuelve camino.''}\\[1mm]
Cada error de hoy, bien observado, se convierte en pieza del camino que pisarás mañana.
\end{softbox}

\small
\begin{tabularx}{\textwidth}{>{\bfseries}p{3.6cm}Y}
\rowcolor{milcGradeDark!12}Nombre completo: & \linead{10cm}\\
Fecha: & \linead{4cm} \quad Curso/grupo: \linead{4cm}\\
Mi firma: & \linead{9cm}\\
\end{tabularx}
\normalsize

\begin{infoband}{milcGradeDark}
\textbf{Para la próxima sesión.} Llega con tu encuesta aplicada y al menos 15-20 respuestas reales recolectadas. En la sesión 6 vas a aprender a hacer \textbf{análisis cualitativo con IA} de las respuestas abiertas. La encuesta de hoy es la materia prima; el análisis de mañana es donde se vuelve información útil.
\end{infoband}

\end{document}
